% Part: counterfactuals
% Chapter: minimal-change-semantics
% Section: transitivity

\documentclass[../../../include/open-logic-section]{subfiles}

\begin{document}

\olfileid{cnt}{min}{tra}

\olsection{संक्रमकता}

वास्तविक अभिव्यंजनासाठी
साखळी-नियम लागू होतो:
$!A \lif !B, !B
\lif !C \Entails !A \lif !C$.
दुसऱ्या शब्दांत,
वास्तविक अभिव्यंजन
संक्रमक आहे.
प्रतिवास्तविक विधानांसाठीही
हे सत्य आहे का?
स्टालनाकर यांनी दिलेले
पुढील उदाहरण पाहा.
\begin{quote}
  जे.~एडगर हूव्हर रशियात
  जन्मले असते, तर ते
  कम्युनिस्ट झाले असते.

  जे.~एडगर हूव्हर कम्युनिस्ट
  असते, तर ते देशद्रोही
  झाले असते.

  म्हणून जे.~एडगर हूव्हर
  रशियात जन्मले असते,
  तर ते देशद्रोही
  झाले असते.
\end{quote}
हूव्हर यांचा जन्म
(प्रत्यक्षात ज्या वेळी झाला
त्याच वेळी) अमेरिकेत
न होता रशियात झाला
असता, तर ते
सोव्हिएत संघात वाढले
असते आणि कम्युनिस्ट
झाले असते, असे
समजू या. म्हणून
पहिले गृहीतक सत्य
आहे. त्याचप्रमाणे,
दुसरे गृहीतक
स्वतंत्रपणे पाहिल्यास
सत्य आहे. पण
निष्कर्ष असत्य आहे:
हूव्हर यांचा जन्म
सोव्हिएत संघात
झाला असता, तर
ते बहुधा निष्ठावान
कम्युनिस्ट झाले
असते; आपल्या
देशाशी विश्वासघात
केला नसता.
सत्य-मूल्यांची ही
सहज वाटणारी
नेमणूक स्टालनाकर--लुईस
यांच्या विश्लेषणातूनही
मिळते. हूव्हर यांच्या
जन्मस्थळाखेरीज
दुसरा बदल नाही
असे आपल्या जगाच्या
सर्वाधिक जवळचे
संभाव्य जग म्हणजे
हूव्हर सोव्हिएत
संघाचे चांगले नागरिक
म्हणून वाढतात
असे जग. पहिल्या
गृहीतकाचे आणि
निष्कर्षाचे पूर्वांग
सत्य असलेले हे
सर्वाधिक जवळचे
संभाव्य जग आहे.
तेथे हूव्हर
कम्युनिस्ट पक्षाचे
निष्ठावान सदस्य
असतात, म्हणून
देशद्रोही नसतात.
पण दुसरे गृहीतक
तपासताना वेगळे
जग पाहावे लागते:
हूव्हर कम्युनिस्ट
असलेले सर्वाधिक
जवळचे जग असे
की ते अमेरिकेत
जन्मले, नंतर
कम्युनिस्ट झाले,
आणि त्यामुळे
देशद्रोही ठरले.\footnote{
  अर्थात या उदाहरणाचा
  मुद्दा समजण्यासाठी
  काही सत्ताविषयक
  आणि राजकीय
  गृहीतके मानावी
  लागतात. उदाहरणार्थ,
  हूव्हर यांचा जन्म
  रशियन पालकांपासून
  होणे शक्य होते,
  किंवा १९५० च्या
  दशकातील अमेरिकेतील
  कम्युनिस्ट हे आपल्या
  देशाचे देशद्रोही होते,
  अशी गृहीतके.}

\begin{prob}
  प्रतिवास्तविक विधानांची
  संक्रमकता अपयशी
  ठरते याचे पटण्याजोगे,
  सहज उदाहरण शोधा.
\end{prob}

\begin{ex}\ollabel{ex:trans-counterex}
  गोलक चिन्हार्थमीमांसेत हे
  अनुमान अवैध ठरते;
  म्हणजे
  $p \cif q, q \cif r
  \Entails/ p \cif r$.
  $\mModel{M} = \tuple{W, O, V}$
  हे प्रतिमान पाहा: येथे
  $W = \{w, w_1, w_2\}$,
  $O_w = \{\{w\}, \{w,
  w_1\}, \{w, w_1, w_2\}\}$,
  $V(p) = \{w_2\}$,
  $V(q) = \{w_1, w_2\}$,
  आणि $V(r) = \{w_1\}$.
  $p$ समाविष्ट करणारा
  $S = \{w, w_1, w_2\}$
  हा गोलक आहे आणि
  त्यातील सर्व जगांत
  $p \lif q$ सत्य आहे;
  म्हणून
  $\mSat{M}{p \cif q}[w]$.
  तसेच $q$ समाविष्ट
  करणारा
  $S' = \{w, w_1\}$
  हा गोलक आहे आणि
  त्यातील सर्व जगांत
  $\mSat{M}{q \lif r}$
  सत्य आहे; म्हणून
  $\mSat{M}{q \cif r}[w]$.
  पण $p$ समाविष्ट
  करणाऱ्या
  $\{w, w_1, w_2\}$
  गोलकात~$w_2$
  हे जग आहे,
  जिथे
  $\mSat/{M}{p \lif
    r}[w_2]$.
\end{ex}

\begin{prob}
  \olref[cnt][min][tra]{ex:trans-counterex}
  मधील प्रतिउदाहरणाला
  अनुरूप गोलकांची
  आकृती काढा.
\end{prob}

\begin{prob}
  \olref[cnt][min][tra]{ex:trans-counterex}
  मध्ये जग~$w_2$
  येथे हूव्हर यांचा
  जन्म रशियात झाला
  आहे, ते कम्युनिस्ट
  आहेत आणि देशद्रोही
  नाहीत. जग~$w_1$
  येथे त्यांचा जन्म
  अमेरिकेत झाला
  आहे, ते कम्युनिस्ट
  आहेत आणि देशद्रोही
  आहेत. या प्रतिमानात
  $w_1$ हे~$w_2$
  पेक्षा~$w$ च्या
  अधिक जवळ आहे.
  हे आवश्यक आहे
  का? हूव्हर यांचा
  रशियात जन्म होणे
  हे त्यांच्या
  कम्युनिस्ट असण्यापेक्षा
  अधिक दूरची
  शक्यता आहे असे
  न मानणारे
  प्रतिउदाहरण
  देता येईल का?
\end{prob}

\end{document}
